\unit{Exposé 190 -- Técnica de descenso I}
\sid{190}
\src{299}
\seminar{12\textsuperscript{o} año, 1959/60, n\textsuperscript{o} 190\\Diciembre de 1959}
\paperhead{TÉCNICA DE DESCENSO Y TEOREMAS DE EXISTENCIA EN GEOMETRÍA ALGEBRAICA\\
I. GENERALIDADES. DESCENSO POR MORFISMOS FIELMENTE PLANOS}%
  {por Alexander GROTHENDIECK}

Desde un punto de vista técnico, el presente exposé y los que le seguirán
pueden considerarse variaciones sobre el célebre «teorema 90» de Hilbert.
Parece que A. Weil introdujo el método de descenso en geometría algebraica con
el nombre de «descenso del cuerpo base». Weil se restringe, por lo demás, al
caso de extensiones finitas separables de cuerpos. P. Cartier trató después el
caso de extensiones radiciales de altura 1. Por falta del lenguaje de los
esquemas y, más concretamente, por falta de elementos nilpotentes en los
anillos que se permitía considerar, Cartier no había podido formular la
identidad esencial entre estos dos casos.

En la actualidad, parece que la técnica general de descenso que se expondrá
(unida, cuando proceda, a los teoremas fundamentales de la «geometría formal»,
cf. [3]) está en la base de la mayoría de los teoremas de existencia en
geometría algebraica.\footnote{La fe de erratas de 1962 considera excesiva esta
formulación: es válida en gran medida para las técnicas no proyectivas de los
dos primeros exposés de la serie TDTE I a VI, pero no para las técnicas
proyectivas (TDTE IV, V y VI).} Conviene señalar, además, que esta técnica puede
sin duda trasladarse a la «geometría analítica», y cabe esperar que en un
futuro próximo los especialistas sepan demostrar los análogos «analíticos» de
los teoremas de existencia en geometría formal que se darán en el exposé II.
En cualquier caso, los trabajos recientes de Kodaira--Spencer, cuyos métodos
parecen inadecuados para definir y estudiar «variedades de módulos» en un
entorno de sus puntos singulares, muestran con bastante claridad la necesidad
de métodos más próximos a la teoría de esquemas (que deberán completar, por
supuesto, las técnicas trascendentes).

En el presente exposé I trataremos el caso de descenso más elemental (indicado
en el título). Sin embargo, las aplicaciones de los teoremas 1, 2 y 3 que
siguen son ya muy numerosas. Nos hemos limitado a dar algunas a modo de
ejemplo, sin intentar conferirles aquí la máxima generalidad posible.

Utilizaremos libremente el lenguaje de los esquemas, para el cual remitimos al
exposé ya citado, así como a [2]. Señalemos, no obstante, expresamente que los
preesquemas considerados en el presente exposé no son necesariamente
noetherianos, y los

\src{300}
morfismos no son necesariamente de tipo finito. Así, si $A$ es un anillo local
noetheriano, con completado $\widehat A$, habrá que considerar el anillo no
noetheriano $\widehat A\otimes_A\widehat A$, y los morfismos de esquemas afines
correspondientes a las inclusiones entre los anillos considerados.

\fsection{A}{Preliminares sobre las categorías}

\subsection*{1. Categorías fibradas, datos de descenso, morfismos de
$\mathcal F$-descenso}
\addcontentsline{toc}{subsection}{1. Categorías fibradas y datos de descenso}

\result{a. DEFINICIÓN}{1.1} Una \emph{categoría fibrada} $\mathcal F$ de base
$\mathcal C$ consta de una categoría $\mathcal C$; para cada
$X\in\mathcal C$, de una categoría $\mathcal F_X$; para cada
$\mathcal C$-morfismo $f:X\longrightarrow Y$, de un funtor
$f^*:\mathcal F_Y\longrightarrow\mathcal F_X$, denotado también por
\[
  f^*(\xi)=\xi\times_Y X
\]
para $\xi\in\mathcal F_Y$ (considerando $X$ como un «objeto de
$\mathcal C$ sobre $Y$», es decir, provisto del morfismo $f$); y, por último,
para dos morfismos componibles
$X\xrightarrow{f}Y\xrightarrow{g}Z$, de un isomorfismo de funtores
\[
  c_{f,g}:(gf)^*\longrightarrow f^*g^*,
\]
sujetos estos datos a las condiciones siguientes:

\begin{enumerate}
\item[(i)] $\operatorname{id}_X^*=\operatorname{id}$;
\item[(ii)] $c_{f,g}$ es el isomorfismo identidad si $f$ o $g$ es un
isomorfismo identidad;
\item[(iii)] dados tres morfismos componibles
$X\xrightarrow{f}Y\xrightarrow{g}Z\xrightarrow{h}T$, el siguiente diagrama de
isomorfismos, formado mediante isomorfismos de la forma $c_{u,v}$, es
conmutativo:
\[
\begin{array}{ccc}
  (h(gf))^*=((hg)f)^* & \longrightarrow & f^*(hg)^* \\
  \big\downarrow && \big\downarrow \\
  (gf)^*h^* & \longrightarrow & (f^*g^*)h^*=f^*(g^*h^*).
\end{array}
\]
\end{enumerate}

\result{EJEMPLO}{1} Sea $\mathcal C$ una categoría en la que existen los
productos fibrados. Se define entonces una categoría fibrada $\mathcal F$ de
base $\mathcal C$ tomando $\mathcal F_X=$ categoría de los objetos de
$\mathcal C$ sobre $X$; el funtor
$f^*:\mathcal F_Y\longrightarrow\mathcal F_X$ correspondiente a un morfismo
$f:X\longrightarrow Y$ se define mediante el producto fibrado
$Z\longmapsto Z\times_YX$.

\result{EJEMPLO}{2} Sea $\mathcal C$ la categoría de los preesquemas y, para
$X\in\mathcal C$, sea $\mathcal F_X$ la categoría de los haces
cuasi-coherentes de módulos sobre $X$. Si $f:X\longrightarrow Y$ es un morfismo
de preesquemas, $f^*:\mathcal F_Y\longrightarrow\mathcal F_X$ es el funtor
imagen inversa

\src{301}
de haces de módulos. Se obtiene así una categoría fibrada de base
$\mathcal C$.

\result{b. DEFINICIÓN}{1.2} Un diagrama
\[
  E\xrightarrow{u}E'
  \overset{v_1}{\underset{v_2}{\rightrightarrows}}E''
\]
de aplicaciones de conjuntos se llama \emph{exacto} si $u$ es una biyección de
$E$ sobre la parte de $E'$ formada por los $x'\in E'$ tales que
$v_1(x')=v_2(x')$.

\result{DEFINICIÓN}{1.3} Sea $\mathcal F$ una categoría fibrada de base
$\mathcal C$; consideremos un diagrama
\[
  \tag{+}
  S\xleftarrow{\alpha}S'
  \overset{\beta_1}{\underset{\beta_2}{\rightrightarrows}}S''
\]
de morfismos en $\mathcal C$ tales que
$\alpha\beta_1=\alpha\beta_2$; este diagrama se llama
$\mathcal F$-\emph{exacto} si, para cada par $\xi,\eta$ de elementos de
$\mathcal F_S$, el siguiente diagrama de aplicaciones de conjuntos
\[
  \tag{+'}
  \operatorname{Hom}(\xi,\eta)
  \xrightarrow{\alpha^*}
  \operatorname{Hom}(\alpha^*\xi,\alpha^*\eta)
  \overset{\beta_1^*}{\underset{\beta_2^*}{\rightrightarrows}}
  \operatorname{Hom}(\gamma^*\xi,\gamma^*\eta),
\]
donde $\gamma=\alpha\beta_1=\alpha\beta_2$, es exacto. En este último diagrama
hemos identificado $\beta_i^*\alpha^*$ con
$(\alpha\beta_i)^*=\gamma^*$ mediante $c_{\beta_i,\alpha}$, para simplificar.

\result{DEFINICIÓN}{1.4} Sea $\mathcal F$ una categoría fibrada de base
$\mathcal C$; consideremos dos morfismos $\beta_1,\beta_2$ de $S''$ en $S'$,
en $\mathcal C$. Sea
$\xi'\in\mathcal F_{S'}$. Se llama \emph{dato de pegado} sobre $\xi'$
(relativo al par $(\beta_1,\beta_2)$) a un isomorfismo de
$\beta_1^*(\xi')$ sobre $\beta_2^*(\xi')$. Si $\xi',\eta'\in\mathcal F_{S'}$
están provistos cada uno de un dato de pegado, un morfismo
$u:\xi'\longrightarrow\eta'$ en $\mathcal F_{S'}$ se dice
\emph{compatible con los datos de pegado} si el diagrama siguiente es
conmutativo:
\[
\begin{array}{ccc}
  \beta_1^*(\xi') & \longrightarrow & \beta_2^*(\xi') \\
  \big\downarrow && \big\downarrow \\
  \beta_1^*(\eta') & \longrightarrow & \beta_2^*(\eta').
\end{array}
\]
De este modo, los objetos de $\mathcal F_{S'}$ provistos de un dato de pegado
(relativo a $\beta_1,\beta_2$) forman una categoría. Si
$\alpha:S'\longrightarrow S$ es un morfismo tal que
$\alpha\beta_1=\alpha\beta_2$, entonces, para todo $\xi\in\mathcal F_S$, el
objeto $\xi'=\alpha^*(\xi)$

\src{302}
de $\mathcal F_{S'}$ está provisto de un dato de pegado canónico, puesto que
\[
  \beta_i^*\alpha^*(\xi)\simeq(\alpha\beta_i)^*(\xi)=\gamma^*(\xi),
\]
tomando de nuevo $\gamma=\alpha\beta_1=\alpha\beta_2$; además, si
$u:\xi\longrightarrow\eta$ es un morfismo en $\mathcal F_S$, entonces
\[
  \alpha^*(u):\alpha^*(\xi)\longrightarrow\alpha^*(\eta)
\]
es un morfismo en $\mathcal F_{S'}$ compatible con los datos de pegado
canónicos. Se obtiene así un funtor canónico de la categoría $\mathcal F_S$ a
la categoría de los objetos de $\mathcal F_{S'}$ provistos de un dato de pegado
relativo al par $(\beta_1,\beta_2)$. Dicho esto, la definición 1.3 también
puede expresarse diciendo que el diagrama (+) de dicha definición es
$\mathcal F$-exacto si el funtor anterior es «plenamente fiel», es decir, si
define una equivalencia de la categoría $\mathcal F_S$ con una subcategoría de
la categoría de los objetos de $\mathcal F_{S'}$ provistos de un dato de pegado
relativo a $(\beta_1,\beta_2)$.

\result{DEFINICIÓN}{1.5} Se dice que un dato de pegado sobre
$\xi'\in\mathcal F_{S'}$ es \emph{efectivo} (respecto de $\alpha$) si
$\xi'$, provisto de este dato, es isomorfo a un $\alpha^*(\xi)$, con
$\xi\in\mathcal F_S$.

Cuando el diagrama (+) es $\mathcal F$-exacto, el objeto $\xi$ de la definición
anterior queda, por tanto, determinado salvo isomorfismo único, y la categoría
$\mathcal F_S$ es equivalente a la categoría de los objetos de
$\mathcal F_{S'}$ provistos de un dato de pegado efectivo.

\medskip\noindent
c. El caso más importante es aquel en que
\[
  S''=S'\times_S S',
\]
siendo las $\beta_i$ las dos proyecciones $p_1$ y $p_2$ de
$S'\times_S S'$ sobre sus dos factores (suponemos ahora que en
$\mathcal C$ existen los productos fibrados). Se tienen entonces dos
condiciones necesarias inmediatas para que sea efectivo un dato de pegado
$\varphi:p_1^*(\xi')\longrightarrow p_2^*(\xi')$ sobre un
$\xi'\in\mathcal F_{S'}$. Se tiene
\[
  \tag{i}
  \Delta^*(\varphi)=\operatorname{id}_{\xi'},
\]
donde $\Delta:S'\longrightarrow S'\times_S S'$ designa el morfismo diagonal,
y donde se ha identificado $\Delta^*p_i^*(\xi')$ con
$(p_i\Delta)^*(\xi')=\xi'$. Se tiene
\[
  \tag{ii}
  p_{31}^*(\varphi)=p_{32}^*(\varphi)\,p_{21}^*(\varphi),
\]

\src{303}
donde $p_{ij}$ designa la proyección canónica de
$S'\times_S S'\times_S S'$ sobre el producto parcial de sus factores
$i$-ésimo y $j$-ésimo.

\result{DEFINICIÓN}{1.6} Se llama \emph{dato de descenso} sobre
$\xi'\in\mathcal F_{S'}$, relativo al morfismo
$\alpha:S'\longrightarrow S$, a un dato de pegado sobre $\xi'$ relativo al par
$(p_1,p_2)$ de proyecciones canónicas
$S'\times_S S'\longrightarrow S'$ que satisface las condiciones (i) y (ii)
anteriores.

\result{DEFINICIÓN}{1.7} Un morfismo $\alpha:S'\longrightarrow S$ se llama
\emph{morfismo de $\mathcal F$-descenso} si el diagrama de morfismos
\[
  S\xleftarrow{\alpha}S'
  \overset{p_1}{\underset{p_2}{\rightrightarrows}}S'\times_S S'
\]
es $\mathcal F$-exacto (definición 1.3); se dice que $\alpha$ es un
\emph{morfismo de $\mathcal F$-descenso estricto} si, además, todo dato de
descenso (definición 1.6) sobre un objeto de $\mathcal F_{S'}$ es efectivo.

Esta última condición también puede enunciarse, de modo más gráfico, así:
«Da lo mismo» dar un objeto de $\mathcal F_S$ o un objeto de
$\mathcal F_{S'}$ provisto de un dato de descenso.

Obsérvese que si un morfismo $\alpha:S'\longrightarrow S$ admite una «sección»
$s:S\longrightarrow S'$ (es decir, un morfismo $s$ tal que
$\alpha s=\operatorname{id}_S$), entonces es un morfismo de
$\mathcal F$-descenso estricto.\footnote{La fe de erratas de 1962 precisa que
no es necesario suponer previamente que $\alpha$ sea un morfismo de
$\mathcal F$-descenso.} En efecto, si $\xi'\in\mathcal F_{S'}$ está provisto de
un dato de descenso relativo a $\alpha$, «proviene» de
$\xi=s^*(\xi')$.

\medskip\noindent
d. Las nociones anteriores pueden presentarse de manera más intuitiva
introduciendo, para un objeto parámetro $T$ de $\mathcal C$ sobre $S$, el
conjunto
\[
  \operatorname{Hom}_S(T,S')=S'(T),
\]
cuyos elementos se denotarán por $t,t'$, etc. Dado un objeto
$\xi'\in\mathcal F_{S'}$, a cada $t\in S'(T)$ le corresponde entonces un objeto
$t^*(\xi')$ de $\mathcal F_T$, que también se denotará por $\xi'_t$. Un dato de
pegado sobre $\xi'$ (relativo a $p_1,p_2$) queda definido entonces dando, para
todo $T$ sobre $S$ y todo par de puntos $t,t'\in S'(T)$, un isomorfismo
\[
  \varphi_{t',t}:\xi'_t\longrightarrow\xi'_{t'}
\]
(que satisface condiciones evidentes de funtorialidad al variar $T$). Las
condiciones (i) y (ii) de 1.c se escriben entonces
\[
  \tag{i bis}
  \varphi_{t,t}=\operatorname{id}_{\xi'_t}
  \qquad\text{para todo $T$ y todo $t\in S'(T)$},
\]

\src{304}
\[
  \tag{ii bis}
  \varphi_{t,t''}=\varphi_{t,t'}\varphi_{t',t''}
  \qquad\text{para todo $T$ y todos $t,t',t''\in S'(T)$}.
\]
Se ve, además, que (ii bis), tomando $t=t'=t''$, implica
$\varphi_{t,t}=\varphi_{t,t}^2$; puesto que $\varphi_{t,t}$ es un isomorfismo
por hipótesis, de aquí se obtiene la relación (i bis), que es de hecho una
consecuencia de (ii bis) (por tanto, (i) es consecuencia de (ii)). Pero si ya
no se supone a priori que los $\varphi_{t,t'}$ sean isomorfismos (es decir, que
$\varphi:p_1^*(\xi')\longrightarrow p_2^*(\xi')$ sea un isomorfismo), entonces
(ii bis) no implica necesariamente (i bis); sin embargo, la conjunción de
(ii bis) y (i bis) implica que los $\varphi_{t,t'}$ son isomorfismos (pues se
tendrá
$\varphi_{t,t'}\varphi_{t',t}=\varphi_{t,t}=\operatorname{id}_{\xi'_t}$).

\subsection*{2. Diagramas exactos y epimorfismos estrictos, morfismos de
descenso. Ejemplos}
\addcontentsline{toc}{subsection}{2. Diagramas exactos y epimorfismos estrictos}

\result{a. DEFINICIÓN}{2.1} Sea $\mathcal C$ una categoría. Un diagrama
\[
  T\xrightarrow{\alpha}T'
  \overset{\beta_1}{\underset{\beta_2}{\rightrightarrows}}T''
\]
de morfismos en $\mathcal C$ se llama \emph{exacto} si, para todo
$Z\in\mathcal C$, el diagrama correspondiente de aplicaciones de conjuntos
\[
  \operatorname{Hom}(Z,T)\longrightarrow
  \operatorname{Hom}(Z,T')
  \rightrightarrows\operatorname{Hom}(Z,T'')
\]
es exacto (definición 1.2). Se dice entonces que $(T,\alpha)$ (o, por abuso de
lenguaje, que $T$) es un \emph{núcleo} del par de morfismos
$(\beta_1,\beta_2)$.

Evidentemente, este núcleo está determinado salvo isomorfismo único. Si
$\mathcal C$ es la categoría de los conjuntos, la definición anterior es
compatible con la definición 1.2. Se define de manera dual la exactitud de un
diagrama
\[
  S\xleftarrow{\alpha}S'
  \overset{\beta_1}{\underset{\beta_2}{\rightrightarrows}}S''
\]
de morfismos en $\mathcal C$; se dice entonces que $(S,\alpha)$ es un
\emph{conúcleo} del par de morfismos $(\beta_1,\beta_2)$.

\result{DEFINICIÓN}{2.2} Un morfismo $\alpha:S'\longrightarrow S$ se llama
\emph{epimorfismo estricto} si es un epimorfismo y si, para todo morfismo
$u:S'\longrightarrow Z$, la siguiente condición necesaria es también
suficiente para que $u$ se factorice como
$S'\longrightarrow S\longrightarrow Z$: para todo $S''\in\mathcal C$ y todo
par de morfismos $\beta_1,\beta_2:S''\longrightarrow S'$ tal que
$\alpha\beta_1=\alpha\beta_2$, se tiene también $u\beta_1=u\beta_2$.

Si existe el producto fibrado $S'\times_S S'$, equivale a decir que el diagrama
\[
  S\xleftarrow{\alpha}S'
  \overset{p_1}{\underset{p_2}{\rightrightarrows}}S'\times_S S'
\]

\src{305}
es exacto, es decir, que $S$ es un conúcleo del par $(p_1,p_2)$. En cualquier
caso, un morfismo conúcleo es un epimorfismo estricto. Señalemos también que un
epimorfismo estricto que sea un monomorfismo es un isomorfismo. Dejamos al
lector desarrollar la noción dual de \emph{monomorfismo estricto}.

Para precisar las relaciones entre la noción de morfismo de
$\mathcal F$-descenso y la de epimorfismo estricto, introducimos además las
definiciones siguientes:

\result{DEFINICIÓN}{2.3} Un morfismo $\alpha:S'\longrightarrow S$ se llama
\emph{epimorfismo universal} (resp. \emph{epimorfismo estricto universal}) si,
para todo $T$ sobre $S$, existe el producto fibrado $T'=S'\times_ST$ y la
proyección $T'\longrightarrow T$ es un epimorfismo (resp. un epimorfismo
estricto).

En las categorías muy buenas (como la categoría de los conjuntos, la
categoría de los haces de conjuntos sobre un espacio topológico, las categorías
abelianas, etc.), las cuatro nociones de epiyectividad así introducidas
coinciden; por el contrario, todas ellas son distintas en una categoría como
la de los preesquemas, o la de los preesquemas sobre un preesquema no vacío
dado $S$, incluso si nos limitamos a $S$-esquemas finitos sobre $S$.

\result{DEFINICIÓN}{2.4} Un morfismo $\alpha:S'\longrightarrow S$ se llama
\emph{morfismo de descenso} (resp. \emph{morfismo de descenso estricto}) si es
un morfismo de $\mathcal F$-descenso (resp. un morfismo de
$\mathcal F$-descenso estricto) (cf. definición 1.7), donde $\mathcal F$ denota
la categoría fibrada de base $\mathcal C$ de los objetos de $\mathcal C$ sobre
objetos de $\mathcal C$ (n\textsuperscript{o} 1, ejemplo 1).

\result{PROPOSICIÓN}{2.1} Si en $\mathcal C$ existen los productos finitos y
los productos fibrados (finitos), entonces, en $\mathcal C$, los morfismos de
descenso coinciden con los epimorfismos estrictos universales.

\result{b. EJEMPLOS}{} Sea $\mathcal C$ la categoría de los preesquemas. Sea
$S\in\mathcal C$, y sean $S'$ y $S''$ dos preesquemas finitos sobre $S$, es
decir, correspondientes a haces de álgebras $\mathcal A'$, $\mathcal A''$
sobre $S$ que, como haces de módulos, son cuasi-coherentes y de tipo finito (es
decir, coherentes si $S$ es localmente noetheriano). Sea
$\alpha:S'\longrightarrow S$ el morfismo estructural de $S'$, y sean
$\beta_1,\beta_2$ dos $S$-morfismos de $S''$ en $S'$, definidos por
homomorfismos de álgebras $\mathcal A'\longrightarrow\mathcal A''$, denotados
también por $\beta_1,\beta_2$. Utilizando el hecho de que un morfismo finito es
cerrado (primer teorema de Cohen--Seidenberg), se demuestra fácilmente que el
diagrama
\[
  \tag{+}
  S\xleftarrow{\alpha}S'
  \overset{\beta_1}{\underset{\beta_2}{\rightrightarrows}}S''
\]
en $\mathcal C$ es exacto si y solo si lo es el diagrama de haces

\src{306}
\[
  \mathcal O_S=\mathcal A\longrightarrow\mathcal A'
  \overset{\beta_1}{\underset{\beta_2}{\rightrightarrows}}\mathcal A''
\]
sobre $S$. En particular, si $\alpha:S'\longrightarrow S$ es un morfismo
finito correspondiente a un haz $\mathcal A'$ de álgebras sobre $S$, entonces
$\alpha$ es un epimorfismo estricto si y solo si el diagrama de haces
\[
  \mathcal O_S=\mathcal A\longrightarrow\mathcal A'
  \overset{p_1}{\underset{p_2}{\rightrightarrows}}
  \mathcal A'\otimes_{\mathcal A}\mathcal A'
\]
es exacto (es un epimorfismo si y solo si
$\mathcal A\longrightarrow\mathcal A'$ es inyectivo). Si $S$ es afín de anillo
$A$, y por tanto $S'$ es afín de anillo $A'$ finito sobre $A$, entonces
$S'\longrightarrow S$ es un epimorfismo estricto si y solo si
$A\longrightarrow A'$ es un isomorfismo de $A$ sobre el subanillo de $A'$
formado por los $x'\in A'$ tales que
\[
  1_{A'}\otimes_Ax'-x'\otimes_A1_{A'}=0
\]
(es un epimorfismo si y solo si $A\longrightarrow A'$ es inyectivo). Como ya
hemos señalado, incluso si $S$ es el esquema de un anillo local artiniano, un
morfismo finito $S'\longrightarrow S$ que sea un epimorfismo no es
necesariamente un epimorfismo estricto. Sin embargo, puede demostrarse que, si
$S$ es un preesquema noetheriano, todo morfismo finito
$S'\longrightarrow S$ que sea un epimorfismo es la composición de una sucesión
finita de epimorfismos estrictos (también finitos). Esto muestra, además, que
la composición de dos epimorfismos estrictos no es necesariamente un
epimorfismo estricto.

\result{c}{} Si $(+)$ es un diagrama exacto de morfismos finitos, entonces,
para todo morfismo plano $T\longrightarrow S$ de preesquemas, el diagrama
obtenido de $(+)$ mediante el cambio de base $T\longrightarrow S$ sigue siendo
exacto. De ello se deduce que, si $X,Y$ son dos $S$-preesquemas, siendo $X$
plano sobre $S$, entonces el siguiente diagrama de aplicaciones de conjuntos
(donde $X',Y'$ son las imágenes inversas de $X,Y$ sobre $S'$, y $X'',Y''$ sus
imágenes inversas sobre $S''$)
\[
  \operatorname{Hom}_S(X,Y)\longrightarrow
  \operatorname{Hom}_{S'}(X',Y')
  \rightrightarrows
  \operatorname{Hom}_{S''}(X'',Y'')
\]
es exacto. En particular, si $\mathcal F$ denota la categoría fibrada cuya base
es la categoría $\mathcal C$ de los preesquemas y tal que, para
$X\in\mathcal C$, $\mathcal F_X$ es la categoría de los $X$-preesquemas
planos, entonces el diagrama $(+)$ es $\mathcal F$-exacto. (Este resultado deja
de ser cierto si no se impone la hipótesis de planitud; en particular, un
epimorfismo estricto finito no es necesariamente un morfismo de descenso.) Se
ve asimismo que $(+)$ es $\mathcal F$-exacto si $\mathcal F$ denota la
categoría fibrada para la cual $\mathcal F_X$ es la categoría de los haces
cuasi-coherentes y planos sobre el preesquema $X$ (también aquí la hipótesis

\src{307}
de planitud es esencial). En uno y otro caso, la cuestión de la efectividad de
un dato de pegado (y, más concretamente, de un dato de descenso cuando
$S''=S'\times_SS'$) sobre un objeto plano sobre $S'$ es delicada (y su
respuesta en diversos casos particulares es uno de los objetos principales de
los presentes exposés). El conferenciante ignora si, para todo epimorfismo
estricto finito $S'\longrightarrow S$, todo dato de descenso sobre un haz
cuasi-coherente plano sobre $S'$ es efectivo (incluso suponiendo que $S$ es el
espectro de un anillo local artiniano y limitándose a los haces localmente
libres de rango 1). Más generalmente, sean $A$ un anillo y $A'$ una
$A$-álgebra (todo es conmutativo) tal que el diagrama de aplicaciones
\[
  A\longrightarrow A'\rightrightarrows A'\otimes_AA'
\]
sea exacto, lo que equivale también a que el morfismo
$S'\longrightarrow S$ correspondiente entre los espectros de $A,A'$ sea un
morfismo de $\mathcal F$-descenso, donde $\mathcal F$ es la categoría fibrada
de los haces cuasi-coherentes planos. Sea $M'$ un $A'$-módulo plano provisto de
un dato de descenso a $A$, es decir, de un isomorfismo
\[
  \varphi:M'\otimes_AA'\xrightarrow{\sim}A'\otimes_AM'
\]
de $A'\otimes_AA'$-módulos que satisface las condiciones (i) y (ii) de 1.c
(cuya explicitación en términos de módulos dejamos al lector). ¿Es efectivo
este dato (respecto de la categoría fibrada de los haces cuasi-coherentes
planos)? Sea $M$ el subconjunto de $M'$ formado por los $x'\in M'$ tales que
\[
  \varphi(x'\otimes_A1_{A'})=1_{A'}\otimes_Ax',
\]
que es un sub-$A$-módulo de $M'$. La inyección canónica
$M\longrightarrow M'$ define un homomorfismo de $A'$-módulos
$M\otimes_AA'\longrightarrow M'$. La efectividad de $\varphi$ significa
entonces lo siguiente: $M$ es un $A$-módulo plano y el homomorfismo anterior es
un isomorfismo.

\result{OBSERVACIÓN}{} En las consideraciones anteriores no habíamos impuesto
ninguna hipótesis de planitud a los morfismos del diagrama $(+)$, lo que nos
obligaba, para disponer de una técnica de descenso, a imponer hipótesis de
planitud a los objetos sobre $S,S'$ considerados. En el parágrafo 2 impondremos
una hipótesis de planitud a $\alpha:S'\longrightarrow S$, lo que nos permitirá
disponer de una técnica de descenso para objetos sobre $S,S'$ que ya no están
sometidos a condición alguna de planitud. En todos los casos interviene una
hipótesis de planitud. Esta es una de las principales razones de la importancia
de la noción de planitud en geometría algebraica (¡cuyo papel no podía
manifestarse mientras uno se limitara a cuerpos base, sobre los cuales, en
efecto, todo es plano!).

\src{308}
\subsection*{3. Aplicación a los morfismos étale}
\addcontentsline{toc}{subsection}{3. Aplicación a los morfismos étale}

Sean $A$ un anillo local y $B$ un álgebra local sobre $A$ cuyo ideal maximal
induce el de $A$. Diremos que $B$ es \emph{étale} sobre $A$ (en lugar de
«no ramificada», expresión usada en otros lugares) si satisface las condiciones
siguientes:
\begin{enumerate}
\item[(i)] $B$ es plana sobre $A$;
\item[(ii)] $B/\mathfrak mB$ es una extensión finita separable de
$A/\mathfrak m=k$ (donde $\mathfrak m$ denota el ideal maximal de $A$).
\end{enumerate}
Cuando $A$ y $B$ son noetherianos y $k$ es algebraicamente cerrado, esto
significa que el homomorfismo $\widehat A\longrightarrow\widehat B$ entre los
completados que prolonga $A\longrightarrow B$ es un isomorfismo. Un morfismo de
tipo finito $f:T\longrightarrow S$ se dice \emph{étale en $x\in T$}, o también
que $T$ es \emph{étale sobre $S$ en $x$}, si $\mathcal O_x$ es étale sobre
$\mathcal O_{f(x)}$; se dice que $f$ es \emph{étale}, o se llama a $f$ un
\emph{morfismo étale}, o se dice que $T$ es \emph{étale sobre $S$}, si $f$ es
étale en todo $x\in T$. Señalemos, además, que si $S$ es localmente
noetheriano, el conjunto de puntos de $T$ en los que $f$ es étale es abierto, y
el uso del «main theorem» de Zariski permite precisar la estructura de $T/S$
en un entorno de tal punto (mediante una ecuación de tipo bien conocido).

Cuando $S$ es un esquema de tipo finito sobre el cuerpo de los números
complejos, le corresponde un espacio analítico $\overline S$ en el sentido de
Serre [5], salvo que $\overline S$ puede tener elementos nilpotentes en su haz
estructural, lo que no cambia nada esencial en [5]. Se ve entonces fácilmente
que $f$ es étale si y solo si
$\overline f:\overline T\longrightarrow\overline S$ lo es, es decir, si todo
punto de $\overline T$ admite un entorno en el que $\overline f$ induce un
isomorfismo sobre un abierto de $\overline S$. En particular, a todo
\emph{recubrimiento étale} $T$ de $S$ (es decir, un morfismo finito étale
$f:T\longrightarrow S$) le corresponde un recubrimiento étale $\overline T$ de
$\overline S$, que es conexo si y solo si $T$ lo es [5]. Se ve además
fácilmente que, si $T,T'$ son dos esquemas étale sobre $S$, entonces la
aplicación natural
\[
  \operatorname{Hom}_S(T,T')\longrightarrow
  \operatorname{Hom}_{\overline S}(\overline T,\overline{T'})
\]
es biyectiva, es decir, el funtor $T\longmapsto\overline T$ de la categoría de
los esquemas étale sobre $S$ a la categoría de los espacios analíticos étale
sobre $\overline S$ es «plenamente fiel» y, por tanto, define una equivalencia
de la primera categoría con una subcategoría de la segunda. Un teorema de
Grauert--Remmert [2] implica que, si $S$ es normal, se obtiene así una
equivalencia entre la categoría de los \emph{recubrimientos étale} de $S$ y la
categoría de los recubrimientos étale (finitos) de $\overline S$, es decir, que
todo recubrimiento étale de $\overline S$ es $\overline S$-isomorfo a un
$\overline T$, donde $T$ es un recubrimiento étale de $S$. Demostremos que el
\emph{teorema de Grauert--Remmert sigue siendo válido sin hipótesis de
normalidad sobre $S$}. Sea, en efecto, primero $S'\longrightarrow S$ un
epimorfismo estricto finito; supongamos demostrado el

\src{309}
teorema para $S'$ y demostremos que será cierto para $S$. Sea
$\mathfrak T$ un recubrimiento étale de $\overline S$ y consideremos su imagen
inversa $\mathfrak T'$ sobre $\overline{S'}$, que corresponde a un haz
analítico coherente $\mathcal A'$ de álgebras sobre $\overline{S'}$, imagen
inversa del haz de álgebras $\mathcal A$ sobre $\overline S$ que define
$\mathfrak T$. Por hipótesis, sobre $\overline{S'}$, $\mathfrak T'$ proviene de
un recubrimiento étale $T'$ de $S'$, es decir, $\mathcal A'$ proviene de un haz
coherente de álgebras $A'$ sobre $S'$. Por otra parte, $\mathcal A'$ está
provisto de un dato de descenso canónico relativo a
$\overline{S'}\longrightarrow\overline S$, es decir, de un isomorfismo entre
sus dos imágenes inversas sobre
\[
  \overline{S'}\times_{\overline S}\overline{S'}
  =\overline{(S'\times_SS')},
\]
(que satisface las condiciones (i), (ii)), y este isomorfismo proviene, según
lo dicho, de un isomorfismo entre los haces algebraicos correspondientes, es
decir, de un dato de descenso sobre $A'$ relativo a
$S'\longrightarrow S$. Se comprueba fácilmente que este último es efectivo
(pues el de $\mathcal A'$ lo es, y la efectividad de un dato de descenso, tal
como se explicitó en el número anterior, se reconoce localmente en los
\emph{completados} de los módulos que intervienen). Se obtiene así un haz
coherente de álgebras $A$ sobre $S$, que define un recubrimiento $T$ de $S$,
el recubrimiento buscado. El resultado anterior sigue siendo entonces
manifiestamente válido si $S'\longrightarrow S$ es solo una composición de un
número finito de epimorfismos estrictos finitos, es decir, si es un
epimorfismo finito cualquiera (por el resultado de factorización señalado en
el parágrafo 2). Se sigue que el teorema de Grauert--Remmert continúa siendo
válido si $S$ es un esquema \emph{reducido}, es decir, si $\mathcal O_S$ no
tiene elementos nilpotentes, como se ve introduciendo su normalización $S'$.
De aquí se pasa fácilmente al caso general.

Una demostración enteramente análoga, que utiliza de nuevo el resultado de
factorización para los epimorfismos estrictos finitos y la naturaleza «formal»
de la efectividad de los datos de descenso, permite demostrar el resultado
siguiente: sea $S$ un preesquema localmente noetheriano y sea
$S'\longrightarrow S$ un morfismo \emph{finito, sobreyectivo y radicial} (o,
lo que equivale a lo mismo, un morfismo de tipo finito tal que, para todo $T$
sobre $S$, el morfismo $T'=S'\times_ST\longrightarrow T$ sea un homeomorfismo,
lo cual también se expresa diciendo que $S'\longrightarrow S$ es un
\emph{homeomorfismo universal}). Para todo $T$ étale sobre $S$, consideremos
su imagen inversa $T'=T\times_SS'$, que es étale sobre $S'$. Entonces el
funtor $T\longmapsto T'$ es una \emph{equivalencia de categorías entre la
categoría de los preesquemas $T$ étale sobre $S$ y la categoría de los
preesquemas $T'$ étale sobre $S'$}. (Se utiliza la biyectividad de
\[
  \operatorname{Hom}_S(T_1,T_2)\longrightarrow
  \operatorname{Hom}_{S'}(T'_1,T'_2)
\]
para dos preesquemas $T_1,T_2$ étale sobre $S$, hecho cuya comprobación directa
es fácil, y el hecho de que el teorema enunciado es cierto si
$S'=(S,\mathcal O_S/\mathcal J)$,

\src{310}
donde $\mathcal J$ es un haz coherente nilpotente de ideales de $\mathcal O_S$
([4], lema 6)). Obsérvese, además, que no suponemos aquí que los $T$
considerados sean finitos sobre $S$. Este resultado implica, en particular,
que el morfismo $S'\longrightarrow S$ induce un isomorfismo del grupo
fundamental $[\,]$ de $S'$ sobre el de $S$ («\emph{invariancia topológica del
grupo fundamental de un preesquema}»).

\subsection*{4. Relaciones con la 1-cohomología}
\addcontentsline{toc}{subsection}{4. Relaciones con la 1-cohomología}

\medskip\noindent
a. Sea $\mathcal C$ una categoría en la que siempre existe el producto de dos
objetos, y sea $T\in\mathcal C$. Para todo conjunto finito
$I\ne\varnothing$, se puede considerar $T^I$; al variar $I$, se obtiene así un
funtor covariante de la categoría de los conjuntos finitos no vacíos a
$\mathcal C$, es decir, lo que puede llamarse un \emph{objeto simplicial} de
$\mathcal C$, denotado por $K_T$. Este último depende covariantemente de $T$;
además, si $u,v$ son dos morfismos $T\longrightarrow T'$, los morfismos
correspondientes $K_T\longrightarrow K_{T'}$ son homotópicos. Diremos que $T$
\emph{domina} a $T'$ si $\operatorname{Hom}(T,T')\ne\varnothing$; esta es una
relación de preorden filtrante creciente en $\mathcal C$. De lo anterior se
deduce que, si $T$ domina a $T'$, existe una clase canónica (salvo homotopía)
de homomorfismos de objetos simpliciales $K_T\longrightarrow K_{T'}$; en
particular, si $K_T$ y $K_{T'}$ son tales que cada uno domina al otro, entonces
$K_T$ y $K_{T'}$ son homotópicamente equivalentes. Sea ahora $F$ un funtor
(contravariante, para fijar ideas) de $\mathcal C$ a una categoría abeliana
$\mathcal C'$; entonces
\[
  C^*(T,F)=F(K_T)
\]
es un objeto cosimplicial de $\mathcal C'$ y, por tanto, define de la manera
habitual un complejo (de cocadenas) en $\mathcal C'$, del cual puede tomarse la
cohomología
\[
  H^*(T,F)=H^*(C^*(T,F))=H^*(F(K_T))
\]
(se podrá poner una $\mathcal C$ como subíndice de $H^*$ si existe posibilidad
de confusión). Este es un funtor cohomológico en $F$, cuya varianza al variar
$T$ se deduce de lo dicho sobre los $K_T$; más precisamente, para $F$ fijo y
$T$ variable en $\mathcal C$ (preordenada por la relación de dominación), los
$H^*(T,F)$ forman un sistema inductivo de objetos graduados de $\mathcal C'$;
en particular, si $T$ y $T'$ son tales que cada uno domina al otro, entonces
$H^*(T,F)$ y $H^*(T',F)$ son canónicamente isomorfos.

Supongamos que en $\mathcal C$ existen los productos fibrados. Para
$S\in\mathcal C$ fijo, se puede aplicar lo anterior a la categoría
$\mathcal C_S$ de los objetos de $\mathcal C$ sobre $S$; escribiremos
$C^*(T/S,F)$ y $H^*(T/S,F)$ en lugar de $C^*(T,F)$ y $H^*(T,F)$ cuando
queramos precisar que trabajamos en la categoría $\mathcal C_S$; así,

\src{311}
$C^*(T/S,F)$ es un complejo de cocadenas en $\mathcal C'$ que, en dimensión
$n$, es igual a
\[
  F(T\times_ST\times_S\cdots\times_ST)
\]
(donde el paréntesis contiene $n+1$ factores).

Obsérvese que, como de costumbre, puede definirse $H^0(T/S,F)$ sin suponer
abeliana la categoría $\mathcal C'$: es el núcleo (definición 2.1), si existe,
del par de morfismos $F(p_i)$ ($i=1,2$)
\[
  F(T)\rightrightarrows F(T\times_ST)
\]
correspondientes a las dos proyecciones
$p_1,p_2:T\times_ST\rightrightarrows T$. En particular, se tendrá un morfismo
natural (llamado de \emph{aumento})
\[
  F(S)\longrightarrow H^0(T/S,F)
\]
que será un isomorfismo en los casos favorables (en particular, si
$T\longrightarrow S$ es un epimorfismo estricto y $F$ es «exacto por la
izquierda»). De igual modo, cuando $F$ toma valores en la categoría de los
grupos en una categoría $\mathcal C''$, también puede definirse $H^1(T/S,F)$;
cuando $\mathcal C''$ es la categoría de los conjuntos (es decir, $F$ toma
valores en la categoría de los grupos ordinarios, no necesariamente
conmutativos), $H^1(T/S,F)$ es el cociente del subgrupo $Z^1(T/S,F)$ de
$C^1(T/S,F)=F(T\times_ST)$ formado por los $g$ tales que
\[
  F(p_{31})(g)=F(p_{32})(g)\,F(p_{21})(g)
\]
por el grupo de operadores $F(T)$, que actúa sobre $C^1(T/S,F)$ y, en
particular, sobre el subconjunto $Z^1(T/S,F)$ mediante
\[
  \rho(g')\cdot g=F(p_2)(g')\,g\,F(p_1)(g')^{-1}.
\]

\medskip\noindent
b. Sea, por ejemplo, $\mathcal F$ una categoría fibrada de base $\mathcal C$.
Sean $\xi,\eta\in\mathcal F_S$ y, para todo $S'$ sobre $S$, sea
\[
  F_{\xi,\eta}(S')=
  \operatorname{Hom}(\xi\times_SS',\eta\times_SS').
\]
Así, $F_{\xi,\eta}$ es un funtor contravariante de $\mathcal C_S$ a la
categoría de los conjuntos. Dicho esto, \emph{afirmar que el morfismo de
aumento}
\[
  F_{\xi,\eta}(S)\longrightarrow H^0(S'/S,F_{\xi,\eta})
\]
\emph{es un isomorfismo para todo par de elementos
$\xi,\eta\in\mathcal F_S$ significa que
$\alpha:S'\longrightarrow S$ es un morfismo de $\mathcal F$-descenso}
(definición 1.7).

\medskip\noindent
c. Pongamos asimismo, para $\xi\in\mathcal F_S$ y todo objeto $S'$ de
$\mathcal C$ sobre $S$,
\[
  G_\xi(S')=\operatorname{Aut}(\xi\times_SS') ;
\]
se ha definido así un funtor contravariante $G_\xi$ de $\mathcal C_S$ a la
categoría de los

\src{312}
grupos. Dicho esto, se constata que $Z^1(S'/S,G_\xi)$ \emph{se identifica
canónicamente con el conjunto de los datos de descenso sobre
$\xi'=\xi\times_SS'$ relativos a $S'\longrightarrow S$} (definición 1.6), y
$H^1(S'/S,G_\xi)$ \emph{se identifica con el conjunto de las clases (salvo
isomorfismo) de objetos de $\mathcal F_{S'}$ provistos de un dato de descenso
relativo a $\alpha:S'\longrightarrow S$ que, como objetos de
$\mathcal F_{S'}$, son isomorfos a $\xi'=\xi\times_SS'$}. Por tanto, si
$\alpha:S'\longrightarrow S$ es un morfismo de $\mathcal F$-descenso
(cf. (b)), entonces $H^1(S'/S,G_\xi)$ \emph{contiene como subconjunto el
conjunto de las clases (salvo isomorfismo) de objetos $\eta$ de
$\mathcal F_S$ tales que $\eta\times_SS'$ sea isomorfo en
$\mathcal F_{S'}$ a $\xi\times_SS'$; y esta inclusión es una identidad si y
solo si todo dato de descenso sobre $\xi'=\xi\times_SS'$ relativo a
$\alpha:S'\longrightarrow S$ es efectivo}. (Este será el caso, en particular,
si $\alpha:S'\longrightarrow S$ es un morfismo de $\mathcal F$-descenso
estricto.)

\result{OBSERVACIÓN}{} Los complejos de cocadenas del tipo $C^*(T/S,F)$
contienen como casos particulares la mayoría de los complejos estándar
conocidos (cohomología de Čech, cohomología de grupos, etc.) y desempeñan un
papel importante en geometría algebraica, en particular en la «cohomología de
Weil» de los preesquemas.

\medskip\noindent
d. \result{EJEMPLO}{1} Sea $S'$ un objeto sobre $S\in\mathcal C$, y sea
$\Gamma$ un grupo de automorfismos de $S'$ tal que $S'$ sea «formalmente
principal sobre $S$, de grupo $\Gamma$», es decir, tal que el morfismo natural
\[
  \Gamma\times S'\longrightarrow S'\times_SS',
\]
donde $\Gamma\times S'$ denota la suma directa de $\Gamma$ copias de $S'$, sea
un isomorfismo. (Se supone que en $\mathcal C$ existen las sumas directas que
intervienen aquí.) Sea $F$ un funtor contravariante de $\mathcal C$ a la
categoría de los grupos abelianos. Entonces $C^*(S'/S,F)$ es canónicamente
isomorfo al grupo simplicial estándar de cocadenas homogéneas
$C^*(\Gamma,F(S'))$ y, por tanto, $H^*(S'/S,F)$ es canónicamente isomorfo a
$H^*(\Gamma,F(S'))$.

\medskip\noindent
e. \result{EJEMPLO}{2} Sea $\mathcal C$ la categoría de los preesquemas. Se
denota por $G_a$ («grupo aditivo») el funtor contravariante de $\mathcal C$ a
la categoría de los grupos abelianos definido por
\[
  G_a(X)=H^0(X,\mathcal O_X).
\]
Se define de igual modo el funtor $G_m$ («grupo multiplicativo») mediante
\[
  G_m(X)=H^0(X,\mathcal O_X)^*
\]
(grupo de los elementos invertibles del anillo $H^0(X,\mathcal O_X)$), y, más
generalmente, el funtor $\operatorname{Gl}(n)$ («grupo lineal de orden $n$»)
mediante

\src{313}
\[
  \operatorname{Gl}(n)(X)=
  \operatorname{Gl}\bigl(n,H^0(X,\mathcal O_X)\bigr),
\]
que es un funtor de $\mathcal C$ a la categoría de los grupos (no
necesariamente conmutativos si $n>1$; para $n=1$ se recupera $G_m$). Además,
puede interpretarse $\operatorname{Gl}(n)$ como un funtor de automorfismos
(cf. (c)) considerando la categoría fibrada $\mathcal F$ de base $\mathcal C$
tal que, para $X\in\mathcal C$, $\mathcal F_X$ sea la categoría de los haces
localmente libres sobre $X$: en efecto, se tiene
\[
  \operatorname{Gl}(n)(X)=
  \operatorname{Aut}_{\mathcal F_X}(\mathcal O_X^n).
\]
De (b) se sigue que, si $\alpha:S'\longrightarrow S$ es un morfismo de
$\mathcal F$-descenso (cf. parágrafo 2.c),
$H^1(S'/S,\operatorname{Gl}(n))$ contiene el conjunto de las clases (salvo
isomorfismo) de haces localmente libres sobre $S$ cuya imagen inversa sobre
$S'$ es isomorfa a $\mathcal O_{S'}^n$, y esta inclusión es una igualdad si y
solo si todo dato de descenso sobre $\mathcal O_{S'}^n$ (relativo a
$\alpha:S'\longrightarrow S$) es efectivo. Cuando $S$ es el espectro de un
anillo local, esto significa, por tanto,
$H^1(S'/S,\operatorname{Gl}(n))=(e)$, pues todo haz localmente libre sobre $S$
es entonces trivial.

Señalemos la equivalencia de las condiciones siguientes para un morfismo
$\alpha:S'\longrightarrow S$:
\begin{enumerate}
\item[(i)] el homomorfismo de aumento
\[
  H^0(S,\mathcal O_S)=G_a(S)\longrightarrow H^0(S'/S,G_a)
\]
es un isomorfismo;
\item[(ii)] $S'\longrightarrow S$ es un morfismo de $\mathcal F$-descenso
(siendo $\mathcal F$ la categoría fibrada de base $\mathcal C$ considerada
arriba).
\end{enumerate}
Si $S'\longrightarrow S$ es finito, estas condiciones equivalen también a:
\begin{enumerate}
\item[(iii)] $S'\longrightarrow S$ es un epimorfismo estricto
(cf. parágrafo 2.c).
\end{enumerate}

Supongamos ahora que $S=\operatorname{Spec}(A)$,
$S'=\operatorname{Spec}(A')$; entonces
\[
  C^n(S'/S,G_a)=C^n(A'/A,G_a)=\bigotimes_A^{n+1}A',
\]
siendo el operador coborde
$C^n(A'/A,G_a)\longrightarrow C^{n+1}(A'/A,G_a)$ la suma alternada de los
operadores de cara
\[
  \partial_i(x_0\otimes x_1\otimes\cdots\otimes x_n)
  =x_0\otimes\cdots\otimes x_{i-1}\otimes1_{A'}\otimes
   x_i\otimes\cdots\otimes x_n.
\]
Del mismo modo,
\[
  C^n(S'/S,G_m)=C^n(A'/A,G_m)
  \simeq\left(\bigotimes_A^{n+1}A'\right)^*,
\]
y las operaciones simpliciales en $C^*(A'/A,G_m)$ están inducidas por las de
$C^*(S'/S,G_a)$. Se explicitan análogamente las operaciones simpliciales en
$C^*(A'/A,\operatorname{Gl}(n))$. \emph{En todos los casos que conoce el
conferenciante, se tiene}
\[
  H^i(A'/A,G_a)=0\qquad\text{para $i>0$},
\]
\emph{y si $A$ es local, se tiene}
\[
  H^1(A'/A,G_m)=0
  \quad\text{y, más generalmente,}\quad
  H^1(A'/A,\operatorname{Gl}(n))=(e)
\]
(cuando $S'\longrightarrow S$ es un morfismo de $\mathcal F$-descenso, es
decir, cuando el diagrama
$A\longrightarrow A'\rightrightarrows A'\otimes_AA'$ es exacto; compárese con
el parágrafo 2.c). Obsérvese que \emph{el «teorema 90» de Hilbert no es sino
la relación}
\src{314}
\[
  H^1(S'/S,G_m)=0
\]
\emph{cuando $A$ es un cuerpo y $A'$ una extensión de Galois finita de este
(cf. ejemplo 1), y también puede expresarse diciendo que, en el caso
considerado, $S'\longrightarrow S$ es un morfismo de descenso estricto para la
categoría fibrada de los haces localmente libres de rango 1.} Es bajo esta
última forma como conviene generalizar el teorema de Hilbert, variando tanto
las hipótesis sobre el morfismo $S'\longrightarrow S$ como las relativas a los
haces cuasi-coherentes considerados.

Señalemos, por último, la equivalencia de las propiedades siguientes cuando
$A$ es un anillo local artiniano de ideal maximal $\mathfrak m$ y $A'$ una
$A$-álgebra (denotando, para todo entero $k>0$, por $A_k$ (resp. $A'_k$) los
anillos $A/\mathfrak m^{k+1}$ (resp. $A'/\mathfrak m^{k+1}A'$)):
\begin{enumerate}
\item[(i)] $H^1(A'_k/A_k,G_a)=0$ para todo $k$.
\item[(ii)] $H^1(A'_k/A_k,G_m)=0$ para todo $k$.
\item[(iii)] $H^1(A'_k/A_k,\operatorname{Gl}(n))=(e)$ para todo $k$ y todo
$n$.
\end{enumerate}
Si $S'\longrightarrow S$ es un epimorfismo estricto, las condiciones
anteriores implican incluso que es un morfismo de descenso estricto para los
módulos libres (de tipo finito o no) sobre $A'$.

\result{OBSERVACIÓN}{} La definición de los grupos $H^i(S'/S,G_m)$, cuando
$S,S'$ son esquemas de cuerpos $A,A'$, se debe a Amitsur. El grupo
$H^2(S'/S,G_m)$ es particularmente interesante como variante «global» del
grupo de Brauer; para esta variante se podrá consultar [1], capítulo VII.

\fsection{B}{Descenso por morfismos fielmente planos}

\subsection*{1. Enunciado de los teoremas de descenso}
\addcontentsline{toc}{subsection}{1. Enunciado de los teoremas de descenso}

\result{DEFINICIÓN}{1.1} Un morfismo
$\alpha:S'\longrightarrow S$ de preesquemas se dice \emph{plano} si, para todo
$x'\in S'$, $\mathcal O_{x'}$ es un módulo plano sobre el anillo
$\mathcal O_{\alpha(x')}$ (es decir, si
$\mathcal O_{x'}\otimes_{\mathcal O_{\alpha(x')}}M$ es un funtor exacto en el
$\mathcal O_{\alpha(x')}$-módulo $M$). Un morfismo se dice
\emph{fielmente plano} si es plano y sobreyectivo.

Por ejemplo, si $S=\operatorname{Spec}(A)$ y
$S'=\operatorname{Spec}(A')$, entonces $S'$ es plano sobre $S$ si y solo si
$A'$ es un $A$-módulo plano, y $S'$ es fielmente plano sobre $S$ si y solo si
$A'$ es un $A$-módulo fielmente plano (es decir, el funtor
$A'\otimes_A M$ en el $A$-módulo $M$ es exacto y fiel); esto significa también,
en la terminología de Serre [5], que el par $(A,A')$ es plano. Si $S'$ es
fielmente plano sobre $S$, entonces el funtor imagen inversa de haces
cuasi-coherentes sobre $S$ es exacto y fiel; en otras palabras, para que una
sucesión de homomorfismos de haces

\src{315}
cuasi-coherentes sobre $S$ sea exacta, es necesario y suficiente que lo sea su
imagen inversa sobre $S'$ (en particular, para que un homomorfismo de haces
cuasi-coherentes sobre $S$ sea un monomorfismo, resp. un epimorfismo, resp. un
isomorfismo, es necesario y suficiente que lo sea su imagen inversa sobre
$S'$). Esta propiedad sigue siendo cierta si nos restringimos sobre un abierto
cualquiera de $S$ y, bajo esta forma, caracteriza a los morfismos fielmente
planos.

\result{DEFINICIÓN}{1.2} Un morfismo $\alpha:S'\longrightarrow S$ se dice
\emph{cuasicompacto} si la imagen inversa de toda parte abierta cuasicompacta
$U$ de $S$ es cuasicompacta (es decir, unión finita de abiertos afines).

Evidentemente, basta comprobar esta propiedad para los abiertos afines de $S$.
Por ejemplo, un morfismo afín (es decir, tal que la imagen inversa de un abierto
afín sea afín) es cuasicompacto.

La clase de los morfismos planos, resp. fielmente planos, resp. cuasicompactos,
es estable por composición y por «cambio de base», y contiene, por supuesto,
los isomorfismos.

\result{TEOREMA}{1} Sea $\alpha:S'\longrightarrow S$ un morfismo de
preesquemas, fielmente plano y cuasicompacto. Entonces $\alpha$ es un
\emph{morfismo de descenso estricto} (A, definición 1.7) para la categoría
fibrada $\mathcal F$ de los haces cuasi-coherentes (A, parágrafo 1, ejemplo 2).

Este enunciado significa dos cosas:
\begin{enumerate}
\item[(i)] Si $F$ y $G$ son dos haces cuasi-coherentes sobre $S$, y $F'$ y
$G'$ sus imágenes inversas sobre $S'$, entonces el homomorfismo natural
\[
  \operatorname{Hom}(F,G)\longrightarrow\operatorname{Hom}(F',G')
\]
es una biyección del primer miembro sobre el subgrupo del segundo formado por
los homomorfismos $F'\longrightarrow G'$ compatibles con los datos de descenso
canónicos sobre estos haces, es decir, aquellos cuyas imágenes inversas por
las dos proyecciones de $S''=S'\times_SS'$ sobre $S'$ dan un mismo
homomorfismo $F''\longrightarrow G''$.
\item[(ii)] Todo haz cuasi-coherente $F'$ sobre $S'$, provisto de un dato de
descenso relativo al morfismo $\alpha:S'\longrightarrow S$
(A, definición 1.6), es isomorfo (provisto de este dato) a la imagen inversa de
un haz cuasi-coherente $F$ sobre $S$.
\end{enumerate}
Tomando $F=\mathcal O_S$ en (i), se obtiene:

\result{COROLARIO}{1} Sea $G$ un haz cuasi-coherente sobre $S$; sean
$G'$ y $G''$ sus imágenes inversas sobre $S'$ y sobre
$S''=S'\times_SS'$, y sean $p_1,p_2$ las dos proyecciones de $S''$ sobre
$S'$. Entonces el siguiente diagrama de aplicaciones de conjuntos

\src{316}
\[
  \Gamma(G)\xrightarrow{\alpha^*}\Gamma(G')
  \underset{p_2^*}{\overset{p_1^*}{\rightrightarrows}}\Gamma(G'')
\]
es exacto (A, definición 1.(a)).

Por otra parte, la conjunción de (i), (ii) y de la definición 1.1 da el

\result{COROLARIO}{2} Sea $G$ como en el corolario 1. Entonces existe una
correspondencia biunívoca entre los subhaces cuasi-coherentes de $G$ y los
subhaces cuasi-coherentes de $G'$ cuyas imágenes inversas sobre $S''$, por las
dos proyecciones $p_1,p_2$, dan el mismo subhaz de $G''$.

Por supuesto, se tiene un enunciado equivalente en términos de haces
cociente. Como se ha visto (A, parágrafo 4, (e)), el teorema 1 debe considerarse
una generalización del «teorema 90» de Hilbert, e implica como casos
particulares diversas formulaciones en términos de 1-cohomología. Para la
demostración, se reduce fácilmente al caso en que $S=\operatorname{Spec}(A)$,
$S'=\operatorname{Spec}(A')$, y para (i) se reduce fácilmente a demostrar el
corolario 1, es decir, la exactitud del diagrama
\[
  M=A\otimes_A M\longrightarrow A'\otimes_A M
  \rightrightarrows A'\otimes_AA'\otimes_A M
\]
para todo $A$-módulo $M$, lo cual se deduce del lema más general:

\result{LEMA}{1.1} Sea $A'$ una $A$-álgebra fielmente plana. Entonces, para
todo $A$-módulo $M$, el complejo aumentado por $M$
$C^*(A'/A,G_a)\otimes_A M$ (cf. A, parágrafo 4, (e)) es una
\emph{resolución} de $M$.

Basta demostrar que el complejo aumentado deducido del anterior mediante el
cambio de base de $A$ a $A'$ satisface las mismas conclusiones. Esto conduce a
comprobar el enunciado cuando se reemplazan $A$ por $A'$ y $A'$ por
$A'\otimes_AA'$, por lo que nos reduce a un caso en el que existe un
homomorfismo de $A$-álgebras $A'\longrightarrow A$ (o, en términos
geométricos, al caso en que $S'$ sobre $S$ admite una sección). En este caso,
esto se deduce de las generalidades de A, parágrafo 4, (a). Señalemos de paso
el siguiente corolario, que generaliza un enunciado bien conocido en
cohomología de Galois (compárese con A, parágrafo 4, (e)):

\result{COROLARIO}{} Si $A'$ es fielmente plana sobre $A$, se tiene
$H^0(A'/A,G_a)=A$ y
\[
  H^i(A'/A,G_a)=0\qquad\text{para $i\geq1$}.
\]

Para demostrar la parte (ii) del teorema 1, se procede como para (i),
reduciéndose al caso en que $S'$ sobre $S$ admite una sección, donde se deduce
de (i) (cf. A, parágrafo 1, (c)).

Evidentemente, se puede variar ad libitum el teorema 1 y sus corolarios
introduciendo diversas estructuras adicionales en los haces (o sistemas de
haces) cuasi-coherentes considerados. Por ejemplo, dar sobre $S$ un haz
\src{317}
cuasi-coherente de álgebras conmutativas «equivale» a dar sobre $S'$ un haz de
este tipo provisto de un dato de descenso relativo a
$\alpha:S'\longrightarrow S$. Teniendo en cuenta la correspondencia funtorial
entre tales haces cuasi-coherentes sobre $S$ y los preesquemas afines sobre
$S$, se obtiene la segunda afirmación del teorema siguiente:

\result{TEOREMA}{2} Sea $\alpha:S'\longrightarrow S$ como en el teorema 1.
Entonces $\alpha$ es un \emph{morfismo de descenso} (no estricto en general)
(A, definición 2.4), y es un \emph{morfismo de descenso estricto} para la
categoría fibrada de los esquemas afines sobre preesquemas
(A, definición 1.7).

La primera afirmación del teorema significa lo siguiente: sean $X,Y$ dos
preesquemas sobre $S$; sean $X',Y'$ sus imágenes inversas sobre $S'$, y
$X'',Y''$ sus imágenes inversas sobre $S''=S'\times_SS'$. Entonces el siguiente
diagrama de aplicaciones naturales
\[
  \operatorname{Hom}_S(X,Y)\xrightarrow{\alpha^*}
  \operatorname{Hom}_{S'}(X',Y')
  \underset{p_2^*}{\overset{p_1^*}{\rightrightarrows}}
  \operatorname{Hom}_{S''}(X'',Y'')
\]
es exacto; es decir, $\alpha^*$ es una biyección de
$\operatorname{Hom}_S(X,Y)$ sobre la parte de
$\operatorname{Hom}_{S'}(X',Y')$ formada por los homomorfismos compatibles con
los datos de descenso canónicos sobre $X',Y'$ (es decir, aquellos cuyas
imágenes inversas por las dos proyecciones de $S''$ sobre $S'$ son iguales).
Esto se deduce fácilmente del teorema 1, corolario 1, cuando uno se restringe
a $Y$ afín sobre $S$; en el caso general, hay que combinar el teorema 1 con el
resultado siguiente:

\result{LEMA}{1.2} Sea $\alpha:S'\longrightarrow S$ un morfismo fielmente plano
y cuasicompacto. Entonces $S$ se identifica con un \emph{espacio topológico
cociente} de $S'$; es decir, toda parte $U$ de $S$ tal que
$\alpha^{-1}(U)$ sea abierta es abierta.

Para completar el teorema 2, hay que dar criterios de efectividad para un dato
de descenso sobre un $S'$-preesquema $X'$ (cuando no se supone que $X'$ sea
afín sobre $S'$). Señalemos primero que \emph{tal dato de descenso no es
necesariamente efectivo}, incluso si $S$ es el espectro de un cuerpo $k$, $S'$
el espectro de una extensión cuadrática $k'$ de este y $X'$ un esquema
algebraico propio de dimensión 2 sobre $S'$ (como puede verse, según Serre,
utilizando la superficie no proyectiva de Nagata). \emph{Para que un dato de
descenso sobre $X'/S'$ relativo a $\alpha:S'\longrightarrow S$ (fielmente
plano y cuasicompacto) sea efectivo, es necesario y suficiente que $X'$ sea
unión de abiertos $X'_i$, afines sobre $S'$, que sean «estables» por el dato de
descenso sobre $X'$.} Esto se cumple con certeza (cualesquiera que sean $X'/S'$
y el dato de descenso sobre $X'$) si el morfismo
$\alpha:S'\longrightarrow S$ es \emph{radicial} (es decir, inyectivo y con
extensiones residuales que

\src{318}
son radiciales). Puede demostrarse también que sigue cumpliéndose si
$\alpha:S'\longrightarrow S$ es \emph{finito} y toda parte finita de $X'$,
contenida en una fibra de $X'$ sobre $S$, está contenida en un abierto de $X'$
afín sobre $S$ (este es el \emph{criterio de Weil}). En particular, esto se
cumple si $X'/S'$ es cuasiproyectivo, y en ese caso puede demostrarse que el
preesquema «descendido» $X/S$ también es cuasiproyectivo (y proyectivo si
$X'/S'$ lo es). En resumen:

\result{TEOREMA}{3} Sea $\alpha:S'\longrightarrow S$ un morfismo de
preesquemas, fielmente plano y cuasicompacto. Si $\alpha$ es \emph{radicial},
es un \emph{morfismo de descenso estricto}. Si $\alpha$ es finito, es un
morfismo de descenso estricto respecto de la categoría fibrada de los
preesquemas cuasiproyectivos (o proyectivos) sobre preesquemas.

\result{OBSERVACIONES}{} Ignoro si, en la segunda afirmación anterior, la
hipótesis de que $\alpha$ sea un morfismo \emph{finito} es realmente necesaria;
en cualquier caso, se comprueba de manera «formal» que puede sustituirse por la
siguiente hipótesis, aparentemente más débil: \emph{para todo punto de $S$
existe un entorno abierto $U$, un $U'$ finito y fielmente plano sobre $U$, y un
$S$-morfismo de $U'$ en $S'$.} Un caso típico que no entra en el anterior es
aquel en que $S=\operatorname{Spec}(A)$,
$S'=\operatorname{Spec}(\widehat A)$, donde $A$ es un anillo local noetheriano
y $\widehat A$ su completado; o también aquel en que $S'$ es cuasifinito sobre
$S$ (es decir, localmente isomorfo a un abierto de un $S$-esquema finito) y no
finito. En estos dos casos, el conferenciante también ignora la respuesta a la
pregunta siguiente: sea $X$ un $S$-esquema tal que $X'=X\times_SS'$ sea
proyectivo sobre $S'$; ¿es cierto que $X$ es proyectivo sobre
$S$?\footnote{La fe de erratas de 1962 responde negativamente: un morfismo
$S'\longrightarrow S$ cuasifinito, étale y sobreyectivo, o un morfismo
$\operatorname{Spec}(\widehat A)\longrightarrow\operatorname{Spec}(A)$, no es
siempre un morfismo de descenso estricto, incluso si $A$ es el anillo local de
una curva algebraica sobre un cuerpo algebraicamente cerrado $k$ y
$S=\operatorname{Spec}(A)$. Se puede encontrar así un morfismo propio y liso
$f:X\longrightarrow S$, que hace de $X$ un fibrado principal de base $S$ bajo
una curva elíptica $E$, tal que $f':X'\longrightarrow S'$ sea proyectivo, pero
$f$ no lo sea. Por tanto, también es un ejemplo de fibrado principal homogéneo
no isotrivial bajo un esquema abeliano.}

\subsection*{2. Aplicación al descenso de ciertas propiedades de los morfismos}
\addcontentsline{toc}{subsection}{2. Descenso de propiedades de los morfismos}

Sea $P$ una clase de morfismos de preesquemas. Sea
$\alpha:S'\longrightarrow S$ un morfismo de preesquemas, sea
$f:X\longrightarrow Y$ un morfismo de $S$-preesquemas y sea
$f':X'\longrightarrow Y'$ la imagen inversa de $f$ por $\alpha$. Cabe
preguntarse entonces si la relación «$f'\in P$» implica «$f\in P$». Resulta que
la respuesta es afirmativa en muchos casos importantes cuando se supone que
$\alpha$ es \emph{fielmente plano y cuasicompacto} (siendo a veces superflua
esta última hipótesis). Esto se ve directamente y sin dificultad si $P$ es la
clase de los morfismos sobreyectivos, resp. radiciales (estos casos se deducen
de la sobreyectividad de $\alpha$), resp. planos, resp. fielmente planos, resp.
lisos (estos casos se deducen de la planitud fiel de $\alpha$), resp. de tipo
finito. Utilizando los teoremas 1 y 2 y el lema 1.2, se ve también que ocurre
lo mismo si $P$ es una de las clases siguientes: isomorfismos, inmersiones
abiertas, inmersiones cerradas, inmersiones (si $f$ es de tipo finito e $Y$
localmente noetheriano), morfismos afines, morfismos finitos, morfismos
cuasifinitos, morfismos abiertos, morfismos cerrados, homeomorfismos, morfismos
separados,

\src{319}
morfismos propios. El único caso importante no dilucidado es el de los
morfismos proyectivos y cuasiproyectivos, ya señalado en la observación del
parágrafo 1.

\subsection*{3. Descenso por morfismos finitos fielmente planos}
\addcontentsline{toc}{subsection}{3. Descenso por morfismos finitos fielmente planos}

Sea $\alpha:S'\longrightarrow S$ un morfismo finito correspondiente a un haz de
álgebras $\mathcal A'$ sobre $S$ que, como haz de módulos, es localmente libre
de tipo finito y es $\ne0$ en todas partes; entonces $\alpha$ es un morfismo
fielmente plano y cuasicompacto, al que pueden aplicarse los resultados
anteriores. Dar un haz cuasi-coherente $F'$ sobre $S'$ equivale a dar el haz
cuasi-coherente $\alpha_*(F')$ sobre $S$, provisto de su estructura de
$\mathcal A'$-módulo (observando que
$\mathcal A'=\alpha_*(\mathcal O_{S'})$). Para simplificar, este haz sobre $S$
también se denotará por $F'$. Las dos imágenes inversas $p_i^*(F')$ de $F'$
sobre $S'\times_SS'$ corresponden asimismo a los haces cuasi-coherentes de
$(\mathcal A'\otimes_{\mathcal O_S}\mathcal A')$-módulos
\[
  F'\otimes_{\mathcal O_S}\mathcal A'
  \quad\text{y}\quad
  \mathcal A'\otimes_{\mathcal O_S}F'.
\]
Dar un $(S'\times_SS')$-homomorfismo del primero al segundo equivale a dar un
homomorfismo de $(\mathcal A'\otimes_{\mathcal O_S}\mathcal A')$-módulos y,
teniendo en cuenta que $\mathcal A'$ es localmente libre, esto equivale a dar
un homomorfismo de $(\mathcal A'\otimes\mathcal A')$-módulos:
\[
  \mathcal U=
  \operatorname{Hom}_{\mathcal O_S}(\mathcal A',\mathcal A')
  =\mathcal A'\otimes\check{\mathcal A}'
  \longrightarrow
  \operatorname{Hom}_{\mathcal O_S}(F',F'),
\]
es decir, a dar, para cada sección $\xi$ de $\mathcal U$ sobre un abierto $V$,
un homomorfismo de $\mathcal O_S$-módulos
$T_\xi:F'|_V\longrightarrow F'|_V$ que satisface las condiciones
\begin{equation}
  T_{f\xi}(x)=fT_\xi(x),
  \qquad T_{\xi f}(x)=T_\xi(fx),
  \tag{3.1}
\end{equation}
donde $f,x$ son, respectivamente, secciones de $\mathcal A',F'$ sobre un
abierto de $S$ contenido en $V$. Las condiciones (i) y (ii) de un dato de
descenso (A, parágrafo 1, (c)) se escriben entonces, respectivamente:
\begin{equation}
  T_{1_{\mathcal U}}x=x,
  \qquad\text{es decir,}\qquad
  T_{1_{\mathcal U}}=\operatorname{id}_{F'},
  \tag{3.2}
\end{equation}
\begin{equation}
  T_{\xi\eta}=T_\xi T_\eta.
  \tag{3.3}
\end{equation}
En otras palabras, \emph{un dato de descenso sobre $F'$ equivale a una
representación del haz de $\mathcal O_S$-álgebras
$\mathcal U=\operatorname{Hom}_{\mathcal O_S}(\mathcal A',\mathcal A')$ en el
haz de $\mathcal O_S$-álgebras
$\operatorname{Hom}_{\mathcal O_S}(F',F')$ que satisface las dos condiciones de
linealidad (3.1).} Si se tiene un apareamiento de haces cuasi-coherentes sobre
$S'$:
\[
  F'_1\times F'_2\longrightarrow F'_3
\]
(que puede interpretarse como un apareamiento de haces de
$\mathcal A'$-módulos sobre $S$),

\src{320}
y datos de pegado sobre los $F'_i$, definidos por homomorfismos
\[
  T_i:\mathcal U\longrightarrow
  \operatorname{Hom}_{\mathcal O_S}(F'_i,F'_i)
  \qquad (i=1,2,3),
\]
entonces estos datos son compatibles con el apareamiento dado, en el sentido
evidente del término, si y solo si se verifica la condición siguiente.

Para toda sección $\xi$ de $\mathcal U$ sobre un abierto, denotando por
\[
  \Delta\xi=\sum_i\xi'_i\otimes_{\mathcal A'}\xi''_i
\]
la sección de $\mathcal U\otimes_{\mathcal A'}\mathcal U$ (considerando
$\mathcal U$ como $\mathcal A'$-módulo mediante su estructura izquierda)
definida por la fórmula
\[
  \xi\mathbin{\cdot}(fg)=\sum_i\xi'_i(f)\xi''_i(g)
\]
(donde $f$ y $g$ son dos secciones de $\mathcal A'$ sobre un abierto más
pequeño), se tiene la fórmula
\begin{equation}
  T^{(3)}_\xi(x\mathbin{\cdot}y)
  =\sum_iT^{(1)}_{\xi'_i}x\mathbin{\cdot}T^{(2)}_{\xi''_i}y
  \tag{3.4}
\end{equation}
para todo par de secciones $x,y$ de $F'_1,F'_2$, respectivamente, sobre un
abierto más pequeño.\footnote{El texto impreso dice «de $\mathcal A'$». Los
dominios de $T^{(1)}$ y $T^{(2)}$, así como el apareamiento
$F'_1\times F'_2\longrightarrow F'_3$, obligan a que $x$ e $y$ sean secciones
de $F'_1$ y $F'_2$, respectivamente.} (Esta propiedad podrá expresarse
diciendo que los homomorfismos $T^{(i)}$ son compatibles con la aplicación
diagonal de $\mathcal U$ respecto del apareamiento dado.) En particular, las
fórmulas (3.1) a (3.4) permiten interpretar, en términos de representaciones de
álgebras con aplicaciones diagonales, los datos de descenso sobre un haz
cuasi-coherente de álgebras sobre $S'$ y, por tanto, también (restringiéndose a
las álgebras conmutativas) los datos de descenso sobre un $S'$-esquema afín.

De aquí se pasa a una interpretación análoga de los datos de descenso sobre un
$S'$-preesquema $X'$ cualquiera: dar tal $X'$ equivale a dar un preesquema
$X'$ sobre $S$, provisto de un homomorfismo de $\mathcal O_S$-álgebras
\[
  \mathcal A'\longrightarrow\mathcal O_{X'},
\]
y un dato de descenso sobre $X'$ equivale a dar un homomorfismo de haces
(compatible con el morfismo $h:X'\longrightarrow S'$)
\[
  \mathcal U\longrightarrow
  \operatorname{Hom}_{h^{-1}(\mathcal O_S)}
  (\mathcal O_{X'},\mathcal O_{X'}),
\]
que satisface condiciones análogas a las condiciones (3.1) a (3.4) anteriores.

\textsc{Ejemplo 1 («Weil»).} -- Supongamos que $S'/S$ es un recubrimiento étale
de Galois de grupo de Galois $\Gamma$ (cf. A, parágrafos 3 y 4.d). Entonces, un
dato de descenso sobre un haz cuasi-coherente $F'$ sobre $S'$ (resp. sobre un
$S'$-preesquema $X'$) equivale a dar una representación de $\Gamma$ por
automorfismos de $(S',F')$ (resp. de $(S',X')$) compatible con las operaciones
de $\Gamma$ sobre $S'$. Este resultado

\src{321}
es «formal», es decir, se demuestra en términos de categorías, pero desde el
punto de vista de este número se deduce también de la estructura explícita de
$\mathcal U$ (provisto de su estructura de anillo, del homomorfismo de anillos
$\mathcal A'\longrightarrow\mathcal U$ y de la aplicación diagonal), conocida
por completo gracias al resultado siguiente: $\mathcal U$ admite, como
$\mathcal A'$-módulo izquierdo, una base formada por las secciones de
$\mathcal U$ que corresponden a los elementos de $\Gamma$.

\textsc{Ejemplo 2 («Cartier»).} -- Sea $p$ un número primo; supongamos que
$p\mathcal O_S=0$ (es decir, $\mathcal O_S$ es de característica $p$),
$(\mathcal A')^p\subset\mathcal O_S=\mathcal A$ (es decir, $S'/S$ es radicial
de altura 1) y que el haz de álgebras $\mathcal A'$ sobre $\mathcal A$ admite
localmente una $p$-base, es decir, una familia $(x_i)$ de secciones tal que
$\mathcal A'$ está generada como álgebra por los $x_i$, sometidos a las únicas
condiciones $x_i^p=0$. Suponemos que el conjunto de los $i$ es finito, de
cardinal $n$. Sea $\mathfrak g$ el haz de las $\mathcal A$-derivaciones de
$\mathcal A'$; es un haz de $\mathcal A'$-módulos localmente libre de rango
$n$ y, además, un haz de $p$-álgebras de Lie sobre $\mathcal A$ (pero no sobre
$\mathcal A'$) que satisface la condición
\begin{equation}
  [X,fY]=X(f)Y+f[X,Y].
  \tag{3.5}
\end{equation}

\textsc{Lema.} -- $\mathcal U=
\operatorname{Hom}_{\mathcal O_S}(\mathcal A',\mathcal A')$ está generada,
como $\mathcal O_S$-álgebra provista de un homomorfismo de álgebras
$\mathcal A'\longrightarrow\mathcal U$, por el sub-$\mathcal A'$-módulo
izquierdo $\mathfrak g$, con las relaciones adicionales
\begin{equation}
  \left\{
  \begin{aligned}
    Xf-fX&=X(f),\\
    XY-YX&=[X,Y],\\
    X^p&=X^{(p)}.
  \end{aligned}
  \right.
  \tag{3.6}
\end{equation}

Del lema anterior se deduce que un dato de descenso sobre el haz
cuasi-coherente $F'$ sobre $S'$ equivale a dar, para todo $X\in\mathfrak g$, un
$\mathcal O_S$-endomorfismo $\overline X$ de $F'$ que satisface las condiciones
\begin{align}
  \overline{fX}&=f\overline X, \tag{3.7}\\
  \overline X(fx)&=X(f)x+f\overline X(x), \tag{3.8}\\
  \overline{[X,Y]}&=[\overline X,\overline Y], \tag{3.9}\\
  \overline{X^{(p)}}&=\overline X^p. \tag{3.10}
\end{align}
(Esto es lo que podría llamarse una conexión lineal sobre $F'$, sin curvatura
y compatible con la potencia $p$-ésima.) Se explicita de igual modo la noción
de dato de descenso sobre un $S'$-preesquema $X'$; la relación (3.4) se
sustituye aquí por la condición de que los $\overline X$ sean derivaciones de
$\mathcal O_{X'}$. Como el morfismo $S'\longrightarrow S$ es radicial, el
teorema 3 garantiza que todo dato de descenso de este tipo es

\src{322}
efectivo y, por tanto, define un $S$-preesquema $X$.

Obsérvese que no ha sido necesario imponer ninguna hipótesis de planitud, de no
singularidad ni de finitud de ningún tipo sobre $F'$ resp. $X'$.

\subsection*{4. Aplicación a criterios de racionalidad}
\addcontentsline{toc}{subsection}{4. Aplicación a criterios de racionalidad}

Sea $X$ un $S$-preesquema tal que la imagen directa de $\mathcal O_X$ sobre
$S$ sea $\mathcal O_S$; esta propiedad seguirá siendo cierta tras todo cambio
plano $S'\longrightarrow S$ de la base $S$. Si $F$ es un haz invertible (es
decir, localmente libre de rango 1) sobre $X$, los automorfismos de $F$, que se
identifican con las secciones invertibles de $\mathcal O_X$, corresponden
biunívocamente a las secciones invertibles de $\mathcal O_S$. Sea entonces $s$
una sección de $X$ sobre $S$; llamaremos, de manera gráfica, sección de $F$
sobre $s$ a una sección del haz invertible $s^*(F)$ sobre $S$. De lo anterior
se deduce que, si $F_i$ ($i=1,2$) son dos haces invertibles sobre $X$, cada uno
provisto de una sección sobre $s$, y si $F_1$ y $F_2$ son isomorfos, existe un
único isomorfismo de $F_1$ sobre $F_2$ compatible con las secciones en cuestión
(es decir, que transforma la primera en la segunda). Por otra parte, e
independientemente de la sección $s$, convengamos en considerar equivalentes
dos haces invertibles $F_1$ y $F_2$ sobre $X$ tales que todo punto de $S$ tenga
un entorno abierto $U$ para el cual las restricciones de $F_1$ y $F_2$ a
$X|U$ sean isomorfas. Entonces todo haz invertible $F$ sobre $X$ es equivalente
a un haz invertible $F_1$ provisto de una sección distinguida sobre $s$ (se
toma $F_1=F\,s^*(F)^{-1}$), y $F_1$ está determinado salvo isomorfismo. En
otras palabras, la clasificación de los haces invertibles sobre $X$ salvo
equivalencia es la misma que la clasificación, salvo isomorfismo, de los haces
invertibles provistos de una sección distinguida.

Como estas propiedades siguen siendo ciertas tras un cambio plano
$\alpha:S'\longrightarrow S$ de la base (reemplazando la sección $s$ por su
imagen inversa $s'$ por $\alpha$), se concluye, teniendo en cuenta el teorema 1:

\emph{Sea el preesquema $X/S$ como arriba y supongamos que admite una sección
$s$; sea $\alpha:S'\longrightarrow S$ un morfismo fielmente plano y
cuasicompacto; sea $F'$ un haz invertible sobre $X'=X\times_SS'$. Para que $F'$
sea equivalente a la imagen inversa sobre $X'$ de un haz invertible $F$ sobre
$X$, es necesario y suficiente que sus imágenes inversas $p_1^*(F')$ y
$p_2^*(F')$ sobre $X'\times_XX'=X\times_S(S'\times_SS')$ sean equivalentes. Si
es así, $F$ está determinado salvo equivalencia.} (Se dirá entonces que $F'$
es racional sobre $S$.)

Inspirándose en este principio cuando $\alpha:S'\longrightarrow S$ es como en
el ejemplo 1 o el ejemplo 2 del número anterior, se obtienen los criterios de
racionalidad de Weil o de Cartier. (Obsérvese que estos autores se limitan al
caso en que $S$ y $S'$

\src{323}
son espectros de cuerpos; a fortiori, $S$ es entonces el espectro de un anillo
local, y la relación de equivalencia introducida arriba no es sino la relación
de isomorfía.) En el primer caso, $F'$ es racional sobre $S$ si y solo si sus
transformados por $\Gamma$ son equivalentes a $F'$. Para expresar el criterio
de racionalidad en el segundo caso, se consideran, en general, el morfismo
diagonal
\[
  X'\longrightarrow X''=X'\times_XX'
\]
de $X'/X$, el haz de ideales $\mathcal I$ correspondiente sobre
$X'\times_XX'$ y el haz $\mathcal I/\mathcal I^2$, que se identifica con su
imagen inversa $\Omega^1_{X'/X}$ sobre $X'$. Como las restricciones de los
$F''_i=p_i^*(F')$ ($i=1,2$) a la diagonal son isomorfas (pues son isomorfas a
$F'$), es decir, $F''_1(F''_2)^{-1}=F''$ tiene una restricción a la diagonal
que es trivial, se sigue que la restricción de $F''$ a
$(X'',\mathcal O_{X''}/\mathcal I^2)$ viene dada, salvo isomorfismo, por un
elemento bien determinado $\xi$ de
\[
  H^1(X'',\mathcal I/\mathcal I^2)
  =H^1(X',\Omega^1_{X'/X}).
\]
Además, en este caso se tiene
\[
  \Omega^1_{X'/X}
  =\Omega^1_{S'/S}\otimes_{\mathcal O_S}\mathcal O_X,
\]
y, por tanto, si $\Omega^1_{S'/S}$ es localmente libre sobre $S$ (como en el
caso de Cartier), $\xi$ define una sección de
\[
  R^1f'_*(\mathcal O_{X'})\otimes\Omega^1_{S'/S}
\]
sobre $S'$ (llamada clase de Atiyah--Cartier del haz invertible $F'$ sobre
$X'/S$), cuya anulación es necesaria y suficiente para que sean equivalentes
las imágenes inversas de $F'$ por las dos proyecciones de
\[
  (X'',\mathcal O_{X''}/\mathcal I^2)
  =X\times_S(S'',\mathcal O_{S''}/\mathcal J^2)
\]
sobre $X'$ (donde $\mathcal J$ es el haz de ideales sobre
$S''=S'\times_SS'$ definido por el morfismo diagonal
$S'\longrightarrow S'\times_SS'$). Esta anulación es, pues, trivialmente
necesaria para que sean equivalentes las imágenes inversas de $F'$ sobre el
propio $X''=X\times_SS''$ y, por tanto, también para que $F'$ sea equivalente
a la imagen inversa de un haz invertible $F$ sobre $X$. Además, la clase de
Atiyah--Cartier puede interpretarse también como la obstrucción a la existencia,
localmente sobre $S'$, de una conexión de $F'$ relativa a las derivaciones de
$X'/X$; tal conexión queda, además, determinada cuando se conocen las
derivaciones de $F'$ correspondientes a las prolongaciones naturales a $X'$ de
las derivaciones de $S'/S$. De esto y de los desarrollos del número anterior se
concluye fácilmente que, en el caso del ejemplo 2 de dicho número y cuando
$X/S$ admite una sección, la anulación de la clase de Atiyah--Cartier es también
suficiente para que $F'$ sea racional sobre $S$.

\subsection*{5. Aplicación a la restricción de la base en un esquema abeliano}
\addcontentsline{toc}{subsection}{5. Restricción de la base en un esquema abeliano}

Sea $S$ un preesquema. Se llama esquema abeliano sobre $S$ a un esquema $X$
liso y propio

\src{324}
sobre $S$ cuyas fibras en los puntos $x\in S$ son esquemas de variedades
abelianas sobre los $\kappa(x)$. Supongamos que $S$ es noetheriano y regular
(es decir, que sus anillos locales son regulares); entonces, utilizando el
teorema de conexidad de Murre [4] (al menos en el caso de «características
iguales», en el que el teorema citado está demostrado actualmente), puede
demostrarse que toda sección racional de $X$ sobre $S$ está definida en todas
partes (es decir, es una sección), lo que generaliza un teorema clásico de
Weil. Se deduce, más generalmente, que si $X'$ es un esquema liso sobre $S$,
entonces toda $S$-aplicación racional de $X'$ en $X$ está definida en todas
partes. De ello se obtiene lo siguiente, que generaliza un resultado de
Chow--Iang:\footnote{El texto impreso dice «Igusa--Lang»; la fe de erratas de
1962 prescribe «Chow--Iang».} Sea $S$ noetheriano y regular, y denote $K$ su
anillo de funciones racionales (producto directo de cuerpos); sea $X$ un
esquema abeliano sobre $K$. Si $X$ es isomorfo a un $K$-esquema de la forma
$X_0\times_S\operatorname{Spec}(K)$, donde $X_0$ es un esquema abeliano sobre
$S$, entonces $X_0$ está determinado salvo isomorfismo único.

Utilizando el resultado de unicidad anterior, se ve que la cuestión de la
restricción de la base en $X$ es local sobre $S$ (y, por tanto, que basta saber
efectuar la restricción a los $\operatorname{Spec}(\mathcal O_x)$, con
$x\in S$). Se ve de igual manera que si
$\alpha:S'\longrightarrow S$ es un morfismo liso de tipo finito, si $K'$ es el
anillo de funciones racionales de $S'$, y si $X\otimes_KK'$ es de la forma
$X'_0\times_{S'}\operatorname{Spec}(K')$, entonces $X'_0$ está provisto de un
dato de descenso canónico relativo a $\alpha$. Teniendo en cuenta el teorema 3,
se concluye:

\result{PROPOSICIÓN}{5.1} Sea $S$ un preesquema noetheriano, regular e
irreducible, de cuerpo de funciones racionales $K$; sea $K'$ una extensión
finita de $K$, no ramificada sobre $S$; sea $S'$ la normalización de $S$ en
$K'$ (que es, por tanto, un recubrimiento étale de $S$); y sea $X$ un esquema
abeliano sobre $K$ tal que $X\otimes_KK'$ sea de la forma
$X'_0\times_{S'}\operatorname{Spec}(K')$, donde $X'_0$ es un esquema abeliano
proyectivo sobre $S'$. Entonces $X$ es de la forma
$X_0\times_S\operatorname{Spec}(K)$, donde $X_0$ es un esquema abeliano
proyectivo sobre $S$.

\result{OBSERVACIONES}{} El conferenciante ignora si puede sustituirse la
hipótesis de que $S'\longrightarrow S$ sea un recubrimiento étale sobreyectivo
(que permite utilizar el teorema 3) por la hipótesis de que sea un morfismo
liso, sobreyectivo y de tipo finito (incluso si se supone que es un morfismo
étale), o si la proposición sigue siendo válida sin suponer que $X'_0$ sea
proyectivo sobre $S'$ (condición que tal vez se cumpla automáticamente).

\subsection*{6. Aplicación a criterios de trivialidad local e isotrivialidad}
\addcontentsline{toc}{subsection}{6. Criterios de trivialidad local e isotrivialidad}

Sean $S$ un preesquema, $G$ un «preesquema en grupos» sobre $S$, y $P$ un
preesquema sobre $S$ sobre el que «actúa $G$» (por la derecha). Se dice que $P$
es formalmente principal homogéneo bajo $G$ si el morfismo bien conocido
\[
  G\times_SP\longrightarrow P\times_SP
\]

\src{325}
inducido por la acción de $G$ sobre $P$ es un isomorfismo. En lo sucesivo,
supondremos que $G$ es plano sobre $S$ (y, por tanto, fielmente plano sobre
$S$), y reservaremos el nombre de fibrado principal homogéneo bajo $G$ para un
fibrado formalmente principal homogéneo $P$ que sea fielmente plano y
cuasicompacto sobre $S$. Es inmediato que esto equivale a decir que puede
encontrarse un cambio fielmente plano y cuasicompacto
$S'\longrightarrow S$ de la base $S$, tal que el fibrado formalmente principal
homogéneo
\[
  P'=P\times_SS'
\]
bajo $G'=G\times_SS'$ sea trivial, es decir, isomorfo a $G'$ (es decir, que
admita una sección); en particular, se podrá tomar $S'=P$. Obsérvese además
que, si $S$ es localmente noetheriano, la hipótesis de planitud fiel sobre $P$
equivale a la hipótesis de que
\[
  \widehat P_s=P\times_S\operatorname{Spec}(\widehat{\mathcal O}_s)
\]
sea fielmente plano sobre $\widehat{\mathcal O}_s$ para todo $s\in S$ (donde
$\widehat{\mathcal O}_s$ denota el completado del anillo local
$\mathcal O_s$), como se deduce del hecho de que $\widehat{\mathcal O}_s$ es
fielmente plano sobre $\mathcal O_s$. Además, si $P$ es de tipo finito sobre
$S$ localmente noetheriano, el conjunto de los puntos $s$ que satisfacen la
condición anterior es constructible; por tanto, si $S$ es un «preesquema de
Jacobson» (por ejemplo, un esquema de tipo finito sobre un cuerpo o, más
generalmente, el espectro de un anillo de Jacobson), basta comprobar la
condición en cuestión en los puntos cerrados de $S$. Esto nos reduce al caso en
que la base es el espectro de un anillo local completo $A$. Cuando
$S=\operatorname{Spec}(A)$ ($A$ un anillo local noetheriano completo) y $P$ es
de tipo finito sobre $S$, la planitud fiel de $P/S$ equivale también a la
existencia de un $S'$ finito y plano sobre $S$ tal que $P'$ sea trivial; y, si
además $G$ es liso sobre $S$, puede suponerse que $S'$ es étale sobre $S$. Por
consiguiente, si además el cuerpo residual de $A$ es algebraicamente cerrado
(«caso geométrico»), $P$ es fielmente plano sobre $A$ si y solo si es trivial.
Así pues, si $S$ es un preesquema algebraico sobre un cuerpo algebraicamente
cerrado y $G$ es liso y de tipo finito sobre $S$, se ve que la condición de
planitud fiel sobre $S$ equivale a la condición de trivialidad analítica (SLF)
de Serre ([6], pp. 1--12).

Pueden introducirse otros tipos, más fuertes, de condiciones sobre $P$ que
tengan la naturaleza de una «trivialidad local». Diremos, en particular, que
$P$ es isotrivial (resp. estrictamente isotrivial) si, para todo $s\in S$,
existen un entorno abierto $U$ de $s$ y un morfismo finito y fielmente plano
(resp. un recubrimiento étale sobreyectivo) $U'\longrightarrow U$ tales que
$P'=P\times_SU'$ sea trivial. (Nuestra terminología difiere de la de Serre [6],
quien llama localmente isotrivial a lo que nosotros llamamos estrictamente
isotrivial.) La isotrivialidad estricta es especialmente útil si $G$ es liso
sobre $S$, pero, en cambio, es una noción inadecuada en los demás casos.

Si $G$ es afín sobre $S$, todo fibrado principal homogéneo $P$ bajo $G$ es afín
por el parágrafo 2, de donde, gracias al teorema 2, resulta la posibilidad de
«descender»

\src{326}
tales fibrados por morfismos fielmente planos y cuasicompactos. Tomemos, en
particular, $G=\operatorname{Gl}(n)_S$, definido por la condición de que el
funtor de $S$-preesquemas
\[
  S'\longmapsto\operatorname{Hom}_S(S',G)
\]
(con valores en la categoría de los grupos) se identifique con el funtor
\[
  \operatorname{Gl}(n)(S')
  =\operatorname{Gl}\bigl(n,H^0(S',\mathcal O_{S'})\bigr)
\]
de A, parágrafo 4.e. Utilizando los hechos (i) de que todo fibrado principal
homogéneo bajo $G$ (resp. todo haz localmente libre de rango $n$ sobre $S$) se
vuelve isomorfo al objeto «trivial» $G$ (resp. $\mathcal O_S^n$) tras un cambio
fielmente plano y cuasicompacto adecuado de $S$; (ii) de que pueden descenderse
los objetos del tipo considerado (fibrados principales homogéneos bajo $G$,
resp. haces localmente libres de rango $n$) por tales morfismos; y, por último,
(iii) de que el grupo de automorfismos del fibrado trivial sobre un $S'/S$ es
funtorialmente isomorfo al grupo de automorfismos del haz localmente libre
trivial de rango $n$ sobre $S'$, se concluye «formalmente» que «da lo mismo»
dar sobre $S$ (o sobre un $S'/S$) un fibrado principal homogéneo de grupo $G$
que dar allí un haz localmente libre de rango $n$. (Más precisamente, se tiene
una equivalencia de categorías fibradas.) Se concluye, en particular:

\result{PROPOSICIÓN}{6.1} Todo fibrado principal homogéneo de grupo
$\operatorname{Gl}(n)_S$ es localmente trivial.

Mediante argumentos conocidos se obtiene el mismo resultado para grupos
estructurales tales como $\operatorname{Sl}(n)_S$, $\operatorname{Sp}(n)_S$ y
productos de tales grupos. Se deduce también que, si $F$ es un subgrupo cerrado
de $G=\operatorname{Gl}(n)_S$, plano sobre $S$, tal que existe el cociente
$G/F$, y si $G$ es un fibrado principal homogéneo isotrivial (resp.
estrictamente isotrivial) sobre $G/F$, de grupo $F\times_S(G/F)$, entonces todo
fibrado principal homogéneo de grupo $F$ es isotrivial (resp. estrictamente
isotrivial). Esto se aplica a todos los «grupos lineales» sobre $S$ utilizados
hasta ahora y, en particular, cuando $G=S\times_k\Gamma$, donde $S$ es un
preesquema sobre el cuerpo $k$ y $\Gamma$ un grupo lineal en el sentido clásico
y, en particular, liso, sobre $k$. Esto resuelve, pues, para tales grupos, una
cuestión de Serre (loc. cit.).

Señalemos asimismo que, para la mayoría de los grupos lisos sobre $S$ conocidos,
sean lineales o no, y en cualquier caso para los de la forma
$S\times_k\Gamma$ anterior, puede demostrarse que todo fibrado principal
homogéneo isotrivial es estrictamente isotrivial. Esto resuelve, en particular,
otra cuestión de Serre (loc. cit., pp. 1--14), dado que un fibrado principal
homogéneo obtenido por un descenso a la Cartier (cf. parágrafo 3, ejemplo 2) es
manifiestamente isotrivial.

\result{OBSERVACIÓN}{} Una de las dificultades esenciales en estas cuestiones
(aparte de la cuestión de la existencia de los esquemas cociente $G/F$) es la
falta de criterios de efectividad para un dato de descenso por un morfismo
fielmente plano no finito.

\src{327}
\begin{center}
  \textbf{BIBLIOGRAFÍA}
\end{center}

\begin{description}
  \item[{[1]}] Dieudonné (J.) y Grothendieck (Alexander). --- \emph{Éléments
    de géométrie algébrique}, de próxima aparición en las Publications
    mathématiques de l'Institut des Hautes Études Scientifiques.
  \item[{[2]}] Grauert (H.) und Remmert (R.). --- \emph{Komplexe Räume}, Math.
    Annalen, t. 136, 1958, pp. 245--318.
  \item[{[3]}] Grothendieck (Alexander). --- \emph{Géométrie formelle et
    géométrie algébrique}. Séminaire Bourbaki, t. 11, 1958/59, no 182.
  \item[{[4]}] Murre (J. P.). --- \emph{On a connectedness theorem for a
    birational transformation at a simple point}, Amer. J. Math., t. 80, 1958,
    pp. 3--15.
  \item[{[5]}] Serre (Jean-Pierre). --- \emph{Géométrie algébrique et géométrie
    analytique}, Ann. Inst. Fourier Grenoble, t. 6, 1955--56, pp. 1--42.
  \item[{[6]}] Serre (Jean-Pierre). --- \emph{Espaces fibrés algébriques},
    Séminaire Chevalley, t. 2, 1958: Anneaux de Chow et applications, no 1.
\end{description}
